\documentclass{article}

\usepackage{geometry}
\geometry{verbose,tmargin=2cm,bmargin=2cm,lmargin=2cm,rmargin=2cm}


\title{\Huge \texttt{balance\_conf.nml} \\ \Large Explanation of the configuration parameters}
\author{Markus Markl}

\begin{document}

\maketitle

\begin{itemize}
	\item \textbf{flre\_path} : (String) Path to the KiLCA flre data of the pre run
	\item \textbf{vac\_path} : (String) Similar as flre\_path
	\item \textbf{btor} : (Double) Toroidal magnetic field value at the magnetic axis, in Gauss. 
	\item \textbf{rtor} : (Double) Minor radius of the torus
	\item \textbf{rmin} : (Double) Minimum radial value for the calculation
	\item \textbf{rmax} : (Double) Maximum radial value, i.e. location of the wall
	\item \textbf{repar} : (Double) Small radius of the separatrix
	\item \textbf{npoimin} : (Integer) Number of points (?)
	\item \textbf{gg\_factor} : Generate grid factor (?)
	\item \textbf{gg\_width} : Generate grid width (?)
	\item \textbf{gg\_r\_res} : Generate grid radial resolution (?)
	\item \textbf{Nstorage} : Number of steps to be reserved in memory. Used for time evolution.
	\item \textbf{tmax\_factor} : (Double) Factor for the maximal time value.
	\item \textbf{antenna\_factor} : (Double) maximal antenna factor for the ramp up of the coil current in the time evolution
	\item \textbf{iboutype} : (Integer) Type of boundary condition
	\item \textbf{iwrite} : (?)
	\item \textbf{eps} : Epsilon (?)
	\item \textbf{dperp} : (?)
	\item \textbf{icoll} : (?)
	\item \textbf{Z\_i} : (Integer) Ion charge number
	\item \textbf{am} : (Integer) Ion mass number
	\item \textbf{rb\_cut\_in} : (?)
	\item \textbf{re\_cut\_in} : (?)
	\item \textbf{rb\_cut\_out} : (?)
	\item \textbf{re\_cut\_out} : (?)
	\item \textbf{write\_formfactors} : (Boolean) Write out form factors. Diagnostics output.
	\item \textbf{flag\_run\_time\_evolution} : (Boolean) Enables the time evolution if true
	\item \textbf{stop\_time\_step} : Stop if time step gets too small
	\item \textbf{path2out} : (String) Path to the output hdf5 file where the balance data should be written to
	\item \textbf{path2inp }: (String) Path to the input hdf5 file that is generated in the pre run
	\item \textbf{timstep\_min} : (Double) Minimal time step value
	\item \textbf{paramscan} : (Boolean) Enables the parameter scan mode, i.e. the changing of input profiles by the multiplication of a factor
	\item \textbf{save\_prof\_time\_step} : (Integer) Save the profiles every nth step
	\item \textbf{diagnostics\_output} : (Boolean) If set true, all of the data is written out. If false, only interesting data is written to hdf5. This saves the most disk space.
	\item \textbf{br\_stopping} : (Boolean) Enables the stopping criterion based on a linear regression of $|B_r^{res}|$, i.e. if its increase is above a certain value, the code is stopped.
	\item \textbf{suppression\_mode} : (Boolean) Enables the suppression of the profile output. Only important data ($D_{e22}^{ql}(r_{res})$, $B_r(r_{res})$) is written out. If enabled, this saves a lot of disk space.
	\item \textbf{debug\_mode} : (Boolean) Enables debugging, i.e. a lot more information is written to the log file.
	\item \textbf{timing\_mode} : (Boolean) Enables the timing mode, i.e. get the time each individual parameter scan needs and also the total time of the run (per mode) (is \textbf{WIP}).
	\item \textbf{ramp\_up\_mode} : (Integer) Enables different ramp up modes, 
	\begin{itemize}
		\item 0: standard ramp-up is used i.e. $antenna\_factor = antenna_\_factor_0 + time^2$
		\item 1: the antenna factor reaches its experimental value at \texttt{t\_max\_ramp\_up}
		\item 3: steady state mode. Goes up to antenna factor max and stays there for $10\cdot t\_max\_ramp\_up$.
	\end{itemize}
	\item \textbf{t\_max\_ramp\_up} : (Double) Time value at which the antenna factor reaches the experimental value. Given in seconds.
	\item \textbf{temperature\_limit} : (Double) Lower electron and ion temperature limit in eV. Only relevant in time evolution. Simulations seem to be sensitive to this value. If it is bad, the calculation will get stuck in a permanent loop. Usually, the value is about $10$ to $50eV$. Best way to check if the code got stuck is checking the size of the output hdf5 file. It should increase at least every minute.
	\item \textbf{antenna\_max\_stopping} : (Double) Stopping criterion that is always on. Code will stop when antenna factor is greater than antenna\_factor\_max\footnote{antenna\_factor\_max is given by the initial antenna\_factor contained in the balance\_conf.nml file.} * antenna\_max\_stopping. Usually, it has a value of about 2 to 5.
	\item \textbf{viscosity\_factor} : (Double) Factor $\alpha$ with which the anomalous viscosity is scaled. The relation is 
	\begin{equation}
		\mu_a = \alpha \cdot D_a
	\end{equation}

\end{itemize}


\end{document}
